\notechapter{N-20221227}{Category Theory Continued and the Question Topology}

\section{More examples of categories}

This note continues the category-theory examples.  Abelian groups with group homomorphisms form \(\mathbf{Ab}\).  Modules over a ring \(A\) form \(\mathbf{Mod}_A\).  Rings with unital ring homomorphisms form \(\mathbf{Rings}\).  Topological spaces with continuous maps form \(\mathbf{Top}\), whose isomorphisms are homeomorphisms.

A partially ordered set \((S,\ge)\) can be regarded as a category: the objects are elements of \(S\), and there is exactly one morphism \(x\to y\) iff \(x\ge y\).  Reflexivity gives identities, transitivity gives composition, and antisymmetry ensures isomorphisms are trivial.

\section{Subcategories and functors}

A subcategory \(\mathcal A\subseteq\mathcal C\) has some objects and some morphisms of \(\mathcal C\), closed under identities and composition.  A functor
\[
  F:\mathcal C\to\mathcal D
\]
sends objects to objects and morphisms to morphisms, preserving identities and composition:
\[
  F(1_A)=1_{F(A)},\qquad F(g\circ f)=F(g)\circ F(f).
\]
The diagram in the original page is a standard functoriality square: the image of a composable path under \(F\) is still composable, and the composite maps to the composite.

\section{Forgetful functors}

A forgetful functor forgets structure.  For example,
\[
  \mathbf{Grp}\to\mathbf{Set}
\]
sends a group to its underlying set and a homomorphism to its underlying function.  Similarly,
\[
  \mathbf{Top}\to\mathbf{Set}
\]
forgets open sets, and \(\mathbf{Vect}_k\to\mathbf{Set}\) forgets vector-space operations.

The note's warning is that a forgetful functor is not innocent: once structure is forgotten, many questions cannot be answered anymore.  A bijection of underlying sets need not be an isomorphism of groups, vector spaces, or topological spaces.

\section{A topology generated by a distance-like question}

The second half moves into topology.  The page considers a collection of subsets and asks what topology it generates.  Given a basis-like family \(\mathcal B\), the topology generated by it consists of arbitrary unions of finite intersections of elements of \(\mathcal B\).

The displayed rational-looking function and its graph on the page are used as a reminder: topology is not only about familiar Euclidean open intervals.  One can impose a topology by declaring certain sets to be basic neighborhoods.


\section{Posets as categories, with an example}

A partially ordered set becomes a category in a very economical way.  There is a unique morphism \(x\to y\) exactly when \(x\le y\).  Identity morphisms come from reflexivity \(x\le x\), and composition comes from transitivity: if \(x\le y\) and \(y\le z\), then \(x\le z\).

This example is important because it shows that a category need not look like a collection of algebraic objects and homomorphisms.  A category can encode logic or order.  Meets and joins in a poset become limits and colimits in the corresponding category.

\section{Functorial thinking}

A functor is not merely a map between categories.  It is a structure-preserving translation.  If \(F:\mathcal C\to\mathcal D\), then for every composable pair
\[
  A\xrightarrow{f}B\xrightarrow{g}C
\]
we require
\[
  F(g\circ f)=F(g)\circ F(f).
\]
This condition says that doing a construction after composing maps gives the same result as first applying the construction to each map and then composing.

Examples clarify the point.  The fundamental group construction is a functor from pointed topological spaces to groups.  Homology is a functor from spaces to graded abelian groups.  The forgetful functor from groups to sets forgets multiplication but still sends homomorphisms to functions.

\section{Generated topology}

The topology part of the note concerns how open sets can be generated from a chosen collection.  If \(\mathcal B\) is intended as a basis, two conditions are expected: every point lies in some basis element, and intersections of basis elements are locally refined by basis elements.  Then every open set is a union of basis elements.

If one starts with a subbasis \(\mathcal S\), the generated topology consists of arbitrary unions of finite intersections of elements of \(\mathcal S\).  This explains the phrase “generated by”: we first force the chosen sets to be open, then close under the axioms of topology.

\section{Why category theory appears beside topology}

Topological constructions are full of universal properties.  Product topology is the coarsest topology making the projection maps continuous.  Quotient topology is the finest topology making the quotient map continuous.  These statements are categorical: they characterize objects by how maps into or out of them behave.

Thus the note's shift from categories to topology is natural.  Topology is not only a study of open sets; it is a study of spaces through continuous maps, and that is exactly categorical thinking.

\section{The larger pattern}

Category theory and topology meet through structure-preserving maps.  A functor preserves categorical structure; a continuous map preserves open-set structure.  A forgetful functor deliberately discards structure.  This is why the note naturally moves from categories to topology: both ask what data must be preserved for two objects to be considered the same.


\paragraph{Expanded synthesis.}
For this chapter, the elementary material should be read as a study of what remains invariant when coordinates change. The earlier sections (`More examples of categories`, `Subcategories and functors`, `Forgetful functors`, `A topology generated by a distance-like question`, `Posets as categories, with an example`) compute with arrays, bases, pivots, determinants, or eigenvectors; the deeper object is a linear map together with the spaces it acts on. A matrix is a representation of that map after bases have been chosen. This is why formulas such as
\[
  A' = P^{-1}AP,\qquad \widetilde A = T^{-1}AS
\]
are not cosmetic: they distinguish an intrinsic morphism from a coordinate description.

The structural hierarchy is as follows. Kernels record what a map forgets, images record what it reaches, cokernels record obstructions to solvability, and quotients record information after an equivalence has been imposed. Determinants act on the top exterior power and therefore measure oriented volume; traces are invariant under cyclic rearrangement and become characters in representation theory; adjoints depend on an inner product and lead to self-adjoint spectral theory.

A useful way to return to the computations is to ask, for every formula, which invariant it is measuring. Row reduction measures rank and kernel; diagonalization measures decomposition into invariant subspaces; Gram--Schmidt measures orthogonal projection; positive definiteness measures ellipsoid geometry; tensor notation measures how components transform. The elementary methods are therefore the finite-dimensional laboratory in which duality, quotienting, functoriality, and spectral decomposition can be seen clearly.
